\section{Metrics and Statistical Methods}\label{sec:metrics}
Per system: \emph{detection} = P(\DIV{} $\mid$ divergent label); \emph{false assurance} = P(\NDO{} $\mid$ divergent label); \emph{abstention} = P(\INS{} or failed run) over the primary denominator of 258; \emph{false alarm} = P(\DIV{} $\mid$ no-difference label); \emph{risk--coverage} with coverage = decided/eligible and risk = errors/decided, where errors are \NDO{} on a divergent label or \DIV{} on a no-difference label. Detection and false assurance are the primary family; the rest are secondary. Failed or missing runs count as abstention (intention-to-treat). Baseline outputs map DIVERGE to \DIV{} and NO\_DIFF to \NDO{}. Labels are generator labels or test-oracle truth and are not adjudicated, so ``false assurance'' means ``no difference reported on a pair labelled divergent'', not a verified miss.

\textbf{Inference.} The inference unit is the project (78 clusters among divergent pairs, 69 among no-difference pairs). Rates are ratios with a project-cluster percentile bootstrap (95\%, $B=10000$, seed 42). Paired comparisons use a project-level sign-flip test, an exact McNemar test at pair level as a sensitivity check, and Holm correction within the pre-registered family (4 comparisons $\times$ 2 primary outcomes). The effect-of-interest rule requires the false-assurance improvement to be at least +10 percentage points with a lower confidence bound above +10. With few independent projects and one dominant cluster, percentile bootstrap intervals can undercover; all $p$-values are descriptive because pairs cluster by project, some labels are synthetic and the data are non-pristine. Power is a design-property simulation, not post-hoc power at the observed effect.

\textbf{Attribution of the false-assurance reduction.} For the pairs that B1 wrongly passes, we cross-tabulate the V7-B verdict. A reduction is ``abstention'' when the recovered pairs become \INS{} rather than \DIV{}. We also report a derived, not executed, gate counterfactual that relabels \INS{} verdicts as \NDO{} (Table~\ref{tab:gate}).

Section~\ref{sec:stats} classifies each finding as supported, descriptive, exploratory, inconclusive or unsupported. Table~\ref{tab:notevaluated} lists what is not evaluated.
\begin{table*}[t]
\centering\scriptsize
\caption{NOT EVALUATED register (Phase 2; sources: \path{reports/phase2/PHASE2_CONSOLIDATED_EXPERIMENTAL_ANALYSIS.md}, section 10, and \path{reports/phase2/K1_K5_READINESS_AND_DECISION.md}).}
\label{tab:notevaluated}
\begin{tabular}{p{0.27\textwidth}p{0.14\textwidth}p{0.5\textwidth}}
\hline
Item & Status & Reason \\
\hline
Confirmatory family and independent generalisation (RQ5) & NOT EVALUATED & no sealed holdout; data non-pristine; K5 not attempted; no independent real project-level run \\
Superiority over any like-for-like comparator; matched-coverage false assurance & NOT EVALUATED & only raw counts; no matched-coverage computation \\
Truth for BugsInPy units, two QuixBugs pairs, six suspect labels & NOT EVALUATED & no defensible label; never relabelled \\
K1 (real-pair feasibility) and K2 under validated isolation & DEFERRED (Phase 3) & neither triggered nor cleared; needs Linux/Docker isolation \\
Linux, Docker and OS-level isolation & NOT VALIDATED & macOS \texttt{sandbox-exec} only; policy flagged below minimum \\
Tamper evidence against a deliberate writer & NOT PROVIDED & unkeyed hash chain; no trust anchor \\
Complete determinism: hash-order sensitivity (E2), \texttt{datetime.now} & NOT EVALUATED & hash seed fixed and never varied; clock not virtualised \\
Published tools (Shadow SE, DiffKemp, Testora, Cleverest, Prefactory, etc.) & NOT COMPARED & not implemented or run; no claim relative to them \\
S3, explorer, module E, declared renames, generator entries & NOT IMPLEMENTED & contribute to \texttt{unsupported\_signature} and abstention \\
Component value of replay, control alternation and thresholds & NOT SHOWN & ablations change no verdict (data have almost no nondeterminism) \\
Calibration of the 0.80 branch floor and repeat counts & NOT EVALUATED & defaults are unvalidated; E5 shows only that the floor relabels S1/S2 \\
Refinement of \texttt{unmodelled\_change} & OPEN & rule is conservative and unrefined (UD-2) \\
Cross-language, cross-platform & NOT SUPPORTED & Python on macOS only \\
\hline
\end{tabular}
\end{table*}

